%----------------------------------------------------------------------------------------
%   Statement of Purpose - Nabil Ahmad (South Korea / HCI & AI Labs)
%   Tailored for KAIST, POSTECH, Yonsei, SKKU, SNU
%----------------------------------------------------------------------------------------
\documentclass[10pt,a4paper]{article}

\usepackage[utf8]{inputenc}
\usepackage[margin=1.5cm,top=1.3cm,bottom=1.3cm]{geometry}
\usepackage{titlesec}
\usepackage{hyperref}
\usepackage{xcolor}

\hypersetup{
    colorlinks=true,
    linkcolor=blue!70!black,
    urlcolor=blue!70!black
}

\titleformat{\section}{\normalsize\bfseries\uppercase}{}{0em}{}[\titlerule]
\titlespacing*{\section}{0pt}{6pt}{3pt}

\setlength{\parindent}{15pt}
\setlength{\parskip}{3pt}

\begin{document}
\pagestyle{empty}

%----------------------------------------------------------------------------------------
%   HEADER
%----------------------------------------------------------------------------------------
\begin{center}
    {\large \textbf{STATEMENT OF PURPOSE}} \\[2pt]
    \textbf{Prospective M.S. Applicant $\cdot$ Human-Computer Interaction \& Applied AI Systems} \\[2pt]
    \small \textbf{Nabil Ahmad} \quad$|$\quad Peshawar, Pakistan \quad$|$\quad +92 304 1237882 \quad$|$\quad \href{mailto:nabilroghani96@gmail.com}{nabilroghani96@gmail.com} \\[1pt]
    Portfolio: \href{https://www.nabilahmad.site}{www.nabilahmad.site} \quad$|$\quad GitHub: \href{https://github.com/nabilroghani}{github.com/nabilroghani}
\end{center}
\vspace{-4pt}

%----------------------------------------------------------------------------------------
%   SECTION 1
%----------------------------------------------------------------------------------------
\section{1. Research Motivation \& Academic Foundations}
The rapid evolution of generative artificial intelligence has unlocked unprecedented computational capability, yet an acute bottleneck persists: modern AI systems often remain disconnected from human cognitive workflows, difficult to scrutinize, and poorly adapted to non-expert users. My overarching research aspiration is to bridge full-stack systems engineering with \textbf{Human-Computer Interaction (HCI)} and \textbf{Human-AI Collaboration} to design intelligent, adaptive interfaces that augment human decision-making rather than complicate it. Completing my Bachelor of Science in Computer Science at the University of Peshawar (CGPA: \textbf{3.32 / 4.00}), I have cultivated a rigorous grounding in data structures, algorithms, operating systems, and web architectures, while pursuing a self-directed research trajectory centered on interactive AI systems and user-centered design.

%----------------------------------------------------------------------------------------
%   SECTION 2
%----------------------------------------------------------------------------------------
\section{2. Engineering \& Human-Centered Systems Experience}
Over the past two years, I have engineered and deployed more than ten production-grade distributed applications, treating engineering rigor as the foundation for experimental research. To investigate how intelligent interfaces can bridge linguistic and cognitive disparities in high-stakes domains, I architected \textbf{MedPulse AI}---a clinical management and patient triage platform. By coupling Google Gemini 2.0 with a MongoDB-grounded Retrieval-Augmented Generation (RAG) pipeline, I developed an interaction paradigm providing structured clinical symptom triage alongside one-click bilingual (English/Urdu) prescription simplifications, empowering patients with low health literacy to understand complex medical regimens.

In parallel, my software engineering role at Apptex Software Solutions for international European clients reinforced my commitment to accessible, robust interaction design. On \textbf{Accurit Accountancy}, I implemented automated accessibility engineering compliant with \textbf{WCAG 2.2 Level AA} standards, verified through automated axe-core accessibility auditing and cross-browser End-to-End testing suites in Playwright. Furthermore, through building real-time collaborative architectures such as \textbf{Penitans} (multi-party tele-consultation platform utilizing WebSockets and WebRTC) and \textbf{BeanVerse} (distributed event-driven marketplace optimized for sub-200ms latency via query indexing), I mastered the mechanics of state synchronization and latency mitigation essential for fluid, zero-perceived-latency interactive user experiences.

%----------------------------------------------------------------------------------------
%   SECTION 3
%----------------------------------------------------------------------------------------
\section{3. Proposed Research Direction \& Laboratory Fit}
South Korea's world-renowned computing research institutions, particularly \textbf{[LAB\_NAME]} led by Prof.~\textbf{[PROFESSOR\_NAME]}, represent the ideal intellectual ecosystem for my graduate studies. I am deeply inspired by the laboratory's pioneering research in \textbf{[RECENT\_PAPER\_OR\_TOPIC]}, specifically concerning how interactive systems can scaffold human capability through generative feedback and mixed-initiative interfaces. During my Master's program, I propose investigating \textbf{Generative, User-Adaptive Interfaces for Domain-Specific AI Assistance}. Specifically, I aim to explore how real-time interaction telemetry and contextual vector memory can dynamically modulate interface complexity, presenting actionable explanations that adapt to a user's cognitive load and domain expertise. My background in building responsive frontend architectures (Next.js 16, React 19, TypeScript) and resilient backend pipelines (Node.js, WebSockets, RAG) equips me to immediately contribute to the laboratory's ongoing projects by rapidly turning theoretical HCI concepts into empirical prototypes ready for user evaluations and benchmarking.

%----------------------------------------------------------------------------------------
%   SECTION 4
%----------------------------------------------------------------------------------------
\section{4. Long-Term Career Objectives}
Following the completion of my M.S. degree, I intend to continue directly into a Ph.D. program in Computer Science / HCI, striving to publish impactful research papers at premier venues such as \textbf{ACM CHI, UIST, CSCW, and IUI}. Ultimately, I aim to pursue an academic and research leadership career developing human-centered computing technologies that democratize access to complex AI tools. With unwavering dedication, practical systems fluency, and a strong work ethic, I am prepared to dedicate 50+ hours weekly to rigorous research and collaborative lab endeavors under your esteemed guidance.

\end{document}

